\documentclass[10pt]{extarticle}
\usepackage[a4paper,margin=1.5cm]{geometry}
\usepackage[T1]{fontenc}
\usepackage{amsmath,amssymb}
\usepackage{graphicx}
\usepackage[export]{adjustbox}
\usepackage{multicol}
\usepackage{enumitem}
\usepackage{booktabs,array,tabularx}
\setlength{\columnsep}{0.7cm}
\setlength{\parindent}{0pt}
\setlength{\parskip}{3pt}
\pagestyle{empty}
\makeatletter
% figure registry (figdefs.tex is generated by the extract step)
\newcommand{\qfigdef}[4]{\@namedef{qf@file@#1}{#2}\@namedef{qf@w@#1}{#3}\@namedef{qf@wide@#1}{#4}}
\newcommand{\qf@check}[1]{\@ifundefined{qf@file@#1}{\PackageError{score2}{Unknown figure id #1}{}}{}}
\newcommand{\qf@img}[4][]{\includegraphics[width=#2\dimexpr\@nameuse{qf@w@#4}pt\relax,max width=#3,max height=0.45\textheight,#1]{\@nameuse{qf@file@#4}}}
% Full-width questions: finalize wraps every question that contains a wide figure in \qfullbegin...\qfullend,
% so the WHOLE question (stem, figures, options) is set across the page in reading order, then the two columns
% resume. The question's head - its number, stem and first figure (\qfullsplit marks the end of the head) -
% is kept on one page, so a number is never left at a page foot with its figure overleaf; the rest may break
% across pages normally, so no big gaps are left. Without a \qfullsplit, the whole question is kept together
% if it is under 45% of a page. Never write these yourself.
\newif\ifqf@full
\newif\ifqf@boxed
\newsavebox{\qf@box}
\newcommand{\qf@open}{\global\qf@fulltrue\par\global\qf@boxedtrue
  \setbox\qf@box\vbox\bgroup\hsize\textwidth\linewidth\textwidth}
\newcommand{\qf@emit}[1]{\par\egroup\global\qf@boxedfalse
  \ifdim\dimexpr\ht\qf@box+\dp\qf@box\relax<#1\textheight\noindent\usebox{\qf@box}\par\else\unvbox\qf@box\fi}
\newcommand{\qfullsplit}{\ifqf@boxed\qf@emit{0.85}\fi}
\newcommand{\qf@close}{\ifqf@boxed\qf@emit{0.45}\fi\global\qf@fullfalse}
\newcommand{\qfullbegin}{\end{multicols}\qf@open}
\newcommand{\qfullend}{\qf@close\begin{multicols}{2}}
\newcommand{\qfullstart}{\qf@open}
\newcommand{\qfulldone}{\qf@close}
% \qfig[scale]{id}: block figure, centred, directly where it is written (size set by the tool).
% Fallback if a wide figure is not inside a full-width question: break out of the columns in place - no floats.
\newcommand{\qfig}[2][1]{\qf@check{#2}%
  \ifqf@full\qf@blk{#1}{#2}\else
  \ifnum\@nameuse{qf@wide@#2}=1 \expandafter\@firstoftwo\else\expandafter\@secondoftwo\fi
  {\end{multicols}\begin{center}\qf@img{#1}{\textwidth}{#2}\end{center}\begin{multicols}{2}}%
  {\qf@blk{#1}{#2}}\fi}
\newcommand{\qf@blk}[2]{\par\smallskip{\centering\qf@img{#1}{\linewidth}{#2}\par}\smallskip}
% \qfigi[scale]{id}: inline figure (inside option lists / tables / minipages); never breaks columns.
\newcommand{\qfigi}[2][1]{\qf@check{#2}\qf@img[valign=c]{#1}{\linewidth}{#2}}
\makeatother
% structure macros - lint reads these to check completeness and figure placement
\newcommand{\papersec}[1]{\par\bigskip{\centering\Large\bfseries PAPER-#1\par}\medskip}
\newcommand{\partsec}[2]{\par\medskip\noindent{\large\bfseries PART-#1: #2}\par\smallskip}
\newcommand{\sect}[1]{\par\smallskip\noindent\textbf{#1}\par}
\newcommand{\qn}[1]{\par\smallskip\noindent\textbf{#1.}~}
\newcommand{\paraq}[2]{\par\smallskip\noindent\textbf{Paragraph for Questions #1 and #2}\par}
\newenvironment{qoptions}{\begin{itemize}[leftmargin=1.8em,itemsep=1pt,topsep=2pt,label={}]}{\end{itemize}}
\newcommand{\opt}[2]{\item[(#1)] #2}
% figure registry (generated by score2_tools.py finalize)
\qfigdef{p1_phy_s1_q02_1}{TNAS - Test 03 (28 May 2023)-img/p1_phy_s1_q02_1.png}{79.3}{0}
\qfigdef{p1_phy_s1_q03_A}{TNAS - Test 03 (28 May 2023)-img/p1_phy_s1_q03_A.png}{126.9}{0}
\qfigdef{p1_phy_s1_q03_B}{TNAS - Test 03 (28 May 2023)-img/p1_phy_s1_q03_B.png}{164.5}{0}
\qfigdef{p1_phy_s1_q03_C}{TNAS - Test 03 (28 May 2023)-img/p1_phy_s1_q03_C.png}{126.9}{0}
\qfigdef{p1_phy_s1_q03_D}{TNAS - Test 03 (28 May 2023)-img/p1_phy_s1_q03_D.png}{126.9}{0}
\qfigdef{p1_phy_s1_q04_1}{TNAS - Test 03 (28 May 2023)-img/p1_phy_s1_q04_1.png}{57.8}{0}
\qfigdef{p1_phy_s1_q08_1}{TNAS - Test 03 (28 May 2023)-img/p1_phy_s1_q08_1.png}{106.3}{0}
\qfigdef{p1_phy_s1_q09_1}{TNAS - Test 03 (28 May 2023)-img/p1_phy_s1_q09_1.png}{92.5}{0}
\qfigdef{p1_phy_s1_q11_1}{TNAS - Test 03 (28 May 2023)-img/p1_phy_s1_q11_1.png}{108.3}{0}
\qfigdef{p1_phy_s1_q12_1}{TNAS - Test 03 (28 May 2023)-img/p1_phy_s1_q12_1.png}{63.7}{0}
\qfigdef{p1_phy_s1_q13_1}{TNAS - Test 03 (28 May 2023)-img/p1_phy_s1_q13_1.png}{111.4}{0}
\qfigdef{p1_phy_s1_q14_1}{TNAS - Test 03 (28 May 2023)-img/p1_phy_s1_q14_1.png}{79.6}{0}
\qfigdef{p1_phy_s1_q15_1}{TNAS - Test 03 (28 May 2023)-img/p1_phy_s1_q15_1.png}{133.5}{0}
\qfigdef{p1_phy_s1_q20_1}{TNAS - Test 03 (28 May 2023)-img/p1_phy_s1_q20_1.png}{91.6}{0}
\qfigdef{p1_phy_s2_q04_1}{TNAS - Test 03 (28 May 2023)-img/p1_phy_s2_q04_1.png}{120.7}{0}
\qfigdef{p1_phy_s2_q05_1}{TNAS - Test 03 (28 May 2023)-img/p1_phy_s2_q05_1.png}{103.1}{0}
\qfigdef{p1_phy_s2_q06_1}{TNAS - Test 03 (28 May 2023)-img/p1_phy_s2_q06_1.png}{74.0}{0}
\qfigdef{p1_phy_s2_q07_1}{TNAS - Test 03 (28 May 2023)-img/p1_phy_s2_q07_1.png}{75.5}{0}
\qfigdef{p1_phy_s2_q08_1}{TNAS - Test 03 (28 May 2023)-img/p1_phy_s2_q08_1.png}{93.9}{0}
\qfigdef{p1_phy_s2_q09_1}{TNAS - Test 03 (28 May 2023)-img/p1_phy_s2_q09_1.png}{103.9}{0}
\qfigdef{p1_chem_s1_q01_A}{TNAS - Test 03 (28 May 2023)-img/p1_chem_s1_q01_A.png}{206.2}{0}
\qfigdef{p1_chem_s1_q01_B}{TNAS - Test 03 (28 May 2023)-img/p1_chem_s1_q01_B.png}{207.6}{0}
\qfigdef{p1_chem_s1_q01_C}{TNAS - Test 03 (28 May 2023)-img/p1_chem_s1_q01_C.png}{205.0}{0}
\qfigdef{p1_chem_s1_q01_D}{TNAS - Test 03 (28 May 2023)-img/p1_chem_s1_q01_D.png}{260.9}{1}
\qfigdef{p1_chem_s1_q07_1}{TNAS - Test 03 (28 May 2023)-img/p1_chem_s1_q07_1.png}{71.5}{0}
\qfigdef{p1_chem_s1_q11_1}{TNAS - Test 03 (28 May 2023)-img/p1_chem_s1_q11_1.png}{109.4}{0}
\qfigdef{p1_chem_s1_q13_A}{TNAS - Test 03 (28 May 2023)-img/p1_chem_s1_q13_A.png}{144.2}{0}
\qfigdef{p1_chem_s1_q13_B}{TNAS - Test 03 (28 May 2023)-img/p1_chem_s1_q13_B.png}{108.5}{0}
\qfigdef{p1_chem_s1_q13_C}{TNAS - Test 03 (28 May 2023)-img/p1_chem_s1_q13_C.png}{141.1}{0}
\qfigdef{p1_chem_s1_q13_D}{TNAS - Test 03 (28 May 2023)-img/p1_chem_s1_q13_D.png}{108.5}{0}
\qfigdef{p1_chem_s1_q14_A}{TNAS - Test 03 (28 May 2023)-img/p1_chem_s1_q14_A.png}{66.1}{0}
\qfigdef{p1_chem_s1_q14_B}{TNAS - Test 03 (28 May 2023)-img/p1_chem_s1_q14_B.png}{102.7}{0}
\qfigdef{p1_chem_s1_q14_C}{TNAS - Test 03 (28 May 2023)-img/p1_chem_s1_q14_C.png}{77.4}{0}
\qfigdef{p1_chem_s1_q14_D}{TNAS - Test 03 (28 May 2023)-img/p1_chem_s1_q14_D.png}{79.3}{0}
\qfigdef{p1_chem_s1_q15_1}{TNAS - Test 03 (28 May 2023)-img/p1_chem_s1_q15_1.png}{109.4}{0}
\qfigdef{p1_chem_s1_q15_A}{TNAS - Test 03 (28 May 2023)-img/p1_chem_s1_q15_A.png}{63.9}{0}
\qfigdef{p1_chem_s1_q15_B}{TNAS - Test 03 (28 May 2023)-img/p1_chem_s1_q15_B.png}{63.9}{0}
\qfigdef{p1_chem_s1_q15_C}{TNAS - Test 03 (28 May 2023)-img/p1_chem_s1_q15_C.png}{55.6}{0}
\qfigdef{p1_chem_s1_q15_D}{TNAS - Test 03 (28 May 2023)-img/p1_chem_s1_q15_D.png}{60.8}{0}
\qfigdef{p1_chem_s1_q16_1}{TNAS - Test 03 (28 May 2023)-img/p1_chem_s1_q16_1.png}{155.5}{0}
\qfigdef{p1_chem_s1_q19_1}{TNAS - Test 03 (28 May 2023)-img/p1_chem_s1_q19_1.png}{87.4}{0}
\qfigdef{p1_chem_s1_q19_A}{TNAS - Test 03 (28 May 2023)-img/p1_chem_s1_q19_A.png}{61.5}{0}
\qfigdef{p1_chem_s1_q19_B}{TNAS - Test 03 (28 May 2023)-img/p1_chem_s1_q19_B.png}{58.3}{0}
\qfigdef{p1_chem_s1_q19_C}{TNAS - Test 03 (28 May 2023)-img/p1_chem_s1_q19_C.png}{61.5}{0}
\qfigdef{p1_chem_s1_q19_D}{TNAS - Test 03 (28 May 2023)-img/p1_chem_s1_q19_D.png}{58.3}{0}
\qfigdef{p1_chem_s2_q03_1}{TNAS - Test 03 (28 May 2023)-img/p1_chem_s2_q03_1.png}{158.5}{0}
\qfigdef{p1_chem_s2_q04_1}{TNAS - Test 03 (28 May 2023)-img/p1_chem_s2_q04_1.png}{124.9}{0}
\qfigdef{p1_chem_s2_q04_2}{TNAS - Test 03 (28 May 2023)-img/p1_chem_s2_q04_2.png}{118.3}{0}
\qfigdef{p1_chem_s2_q04_3}{TNAS - Test 03 (28 May 2023)-img/p1_chem_s2_q04_3.png}{98.0}{0}
\qfigdef{p1_chem_s2_q05_1}{TNAS - Test 03 (28 May 2023)-img/p1_chem_s2_q05_1.png}{169.8}{0}
\qfigdef{p1_chem_s2_q09_1}{TNAS - Test 03 (28 May 2023)-img/p1_chem_s2_q09_1.png}{50.4}{0}
\qfigdef{p1_math_s1_q10_1}{TNAS - Test 03 (28 May 2023)-img/p1_math_s1_q10_1.png}{82.8}{0}
\qfigdef{p2_phy_s1_q01_1}{TNAS - Test 03 (28 May 2023)-img/p2_phy_s1_q01_1.png}{171.4}{0}
\qfigdef{p2_phy_s1_q03_1}{TNAS - Test 03 (28 May 2023)-img/p2_phy_s1_q03_1.png}{92.5}{0}
\qfigdef{p2_phy_s1_q03_2}{TNAS - Test 03 (28 May 2023)-img/p2_phy_s1_q03_2.png}{31.8}{0}
\qfigdef{p2_phy_s1_q03_3}{TNAS - Test 03 (28 May 2023)-img/p2_phy_s1_q03_3.png}{36.2}{0}
\qfigdef{p2_phy_s1_q04_1}{TNAS - Test 03 (28 May 2023)-img/p2_phy_s1_q04_1.png}{157.5}{0}
\qfigdef{p2_phy_s1_q05_1}{TNAS - Test 03 (28 May 2023)-img/p2_phy_s1_q05_1.png}{125.5}{0}
\qfigdef{p2_phy_s1_q06_1}{TNAS - Test 03 (28 May 2023)-img/p2_phy_s1_q06_1.png}{103.5}{0}
\qfigdef{p2_phy_s2_q09_1}{TNAS - Test 03 (28 May 2023)-img/p2_phy_s2_q09_1.png}{157.3}{0}
\qfigdef{p2_phy_s2_q09_2}{TNAS - Test 03 (28 May 2023)-img/p2_phy_s2_q09_2.png}{77.3}{0}
\qfigdef{p2_phy_s2_q09_3}{TNAS - Test 03 (28 May 2023)-img/p2_phy_s2_q09_3.png}{74.7}{0}
\qfigdef{p2_phy_s2_q09_4}{TNAS - Test 03 (28 May 2023)-img/p2_phy_s2_q09_4.png}{103.1}{0}
\qfigdef{p2_phy_s2_q09_5}{TNAS - Test 03 (28 May 2023)-img/p2_phy_s2_q09_5.png}{105.8}{0}
\qfigdef{p2_phy_s3_q03_1}{TNAS - Test 03 (28 May 2023)-img/p2_phy_s3_q03_1.png}{94.2}{0}
\qfigdef{p2_phy_s3_q04_1}{TNAS - Test 03 (28 May 2023)-img/p2_phy_s3_q04_1.png}{76.5}{0}
\qfigdef{p2_phy_s3_q05_1}{TNAS - Test 03 (28 May 2023)-img/p2_phy_s3_q05_1.png}{233.4}{0}
\qfigdef{p2_phy_s3_q06_1}{TNAS - Test 03 (28 May 2023)-img/p2_phy_s3_q06_1.png}{66.1}{0}
\qfigdef{p2_phy_s3_q08_1}{TNAS - Test 03 (28 May 2023)-img/p2_phy_s3_q08_1.png}{213.8}{0}
\qfigdef{p2_chem_s1_q04_1}{TNAS - Test 03 (28 May 2023)-img/p2_chem_s1_q04_1.png}{138.5}{0}
\qfigdef{p2_chem_s1_q04_2}{TNAS - Test 03 (28 May 2023)-img/p2_chem_s1_q04_2.png}{127.3}{0}
\qfigdef{p2_chem_s2_q07_1}{TNAS - Test 03 (28 May 2023)-img/p2_chem_s2_q07_1.png}{71.5}{0}
\qfigdef{p2_chem_s2_q07_2}{TNAS - Test 03 (28 May 2023)-img/p2_chem_s2_q07_2.png}{74.0}{0}
\qfigdef{p2_chem_s2_q07_3}{TNAS - Test 03 (28 May 2023)-img/p2_chem_s2_q07_3.png}{58.7}{0}
\qfigdef{p2_chem_s2_q07_4}{TNAS - Test 03 (28 May 2023)-img/p2_chem_s2_q07_4.png}{64.9}{0}
\qfigdef{p2_chem_s3_q04_1}{TNAS - Test 03 (28 May 2023)-img/p2_chem_s3_q04_1.png}{167.2}{0}
\qfigdef{p2_chem_s3_q05_1}{TNAS - Test 03 (28 May 2023)-img/p2_chem_s3_q05_1.png}{208.6}{0}
\qfigdef{p2_chem_s3_q06_1}{TNAS - Test 03 (28 May 2023)-img/p2_chem_s3_q06_1.png}{115.9}{0}
\qfigdef{p2_chem_s3_q06_2}{TNAS - Test 03 (28 May 2023)-img/p2_chem_s3_q06_2.png}{161.5}{0}
\raggedcolumns
\makeatletter\@beginparpenalty=10000\makeatother
\newcolumntype{L}[1]{>{\raggedright\arraybackslash}p{#1}}
\begin{document}
\begin{center}{\Large\bfseries TNAS --- Test 03 (28 May 2023)}\end{center}
\begin{multicols}{2}
\raggedcolumns
\papersec{1}
\partsec{1}{PHYSICS}
\sect{SECTION-I}
\qn{1} In a uniformly charged sphere of charge $2q$ and radius $R$, two identical particles, each of charge $-q$ are embedded in it. If the system is in static equilibrium, the identical particles should be :-
\begin{qoptions}
\opt{A}{at the ends of a diameter of the sphere.}
\opt{B}{one at the centre, other on the surface of the sphere.}
\opt{C}{at two interior points, separated by a distance $2R$.}
\opt{D}{symmetrically placed about the centre with the separation $R$.}
\end{qoptions}
\qn{2} The currents through $4\Omega$ resistor in the circuit shown below, when the switch $S$ is on and off respectively are :-
\qfig{p1_phy_s1_q02_1}
\begin{qoptions}
\opt{A}{$\left(\dfrac97\mathrm A,3\mathrm A\right)$}
\opt{B}{$\left(\dfrac{10}{7}\mathrm A,2\mathrm A\right)$}
\opt{C}{$\left(\dfrac97\mathrm A,2\mathrm A\right)$}
\opt{D}{$\left(\dfrac{10}{7}\mathrm A,3\mathrm A\right)$}
\end{qoptions}
\qn{3} The electric potential energy of a unit positive charged particle, constrained to move along $x$-axis is given by $U(x)=k_1x^2-k_2x^4$, $k_1$ and $k_2$ being positive constants. The electric field along $x$-axis has the functional form given by :-
\begin{qoptions}
\opt{A}{\qfigi{p1_phy_s1_q03_A}}
\opt{B}{\qfigi{p1_phy_s1_q03_B}}
\opt{C}{\qfigi{p1_phy_s1_q03_C}}
\opt{D}{\qfigi{p1_phy_s1_q03_D}}
\end{qoptions}
\qn{4} A solid cylinder of mass $M$ and radius $r$ starts falling freely under gravity at time $t=0$. The tension in both the strings is $T$ and the linear acceleration of the cylinder is $a$. The instantaneous power $P$ developed at a time $t$ by the force of gravity is :-
\qfig{p1_phy_s1_q04_1}
\begin{qoptions}
\opt{A}{$\dfrac{Mg}{6}$}
\opt{B}{$\dfrac{2Mg}{3}$}
\opt{C}{$\dfrac23Mg^2t$}
\opt{D}{$2Mgt$}
\end{qoptions}
\qn{5} The time period of rotation of the earth is 24 hrs. If the earth increases its speed of rotation, at nearly what time period will the weight of body at the equator become zero ? Given $R=6400$ km and $g=10\ \mathrm{ms}^{-2}$
\begin{qoptions}
\opt{A}{$0.8$ hr}
\opt{B}{$1$ hr}
\opt{C}{$72$ min}
\opt{D}{$84$ min}
\end{qoptions}
\qn{6} Three long straight wires parallel to each other are connected in parallel across a battery of negligible internal resistance. The ratio of their resistances are $3:4:5$. What is the ratio of distance of the middle wire from the others if the net force experienced by it is zero ?
\begin{qoptions}
\opt{A}{$4:3$}
\opt{B}{$3:1$}
\opt{C}{$5:3$}
\opt{D}{$2:3$}
\end{qoptions}
\qn{7} The current in a conductor of 40 ohm resistance increases linearly from 5A to 25 in 10 s. Heat liberated in it during this time is nearly :-
\begin{qoptions}
\opt{A}{$10^5$ J}
\opt{B}{$10^4$ J}
\opt{C}{$10^3$ J}
\opt{D}{$100$ J}
\end{qoptions}
\qn{8} A mysterious circuit element and a $5\ \Omega$ resistor are connected in series with a source of emf of negligible internal resistance through a switch. The element behaves like an open switch initially but allows constant current after long time. After closing the switch, a graph of voltage across the element versus time is drawn as shown in the Figure. The maximum current through the element then is :-
\qfig{p1_phy_s1_q08_1}
\begin{qoptions}
\opt{A}{$5$A}
\opt{B}{$2$A}
\opt{C}{$3$A}
\opt{D}{$4$A}
\end{qoptions}
\qn{9} In the given figure, $i=i_1$ when the switch is closed and $i=i_2$ when the switch is open. What is the value of $(i_1+i_2)$ ?
\qfig{p1_phy_s1_q09_1}
\begin{qoptions}
\opt{A}{$3$A}
\opt{B}{$0$A}
\opt{C}{$12$ A}
\opt{D}{$6$ A}
\end{qoptions}
\qn{10} A solid sphere of radius $R_1$ and volume charge density $\rho=\dfrac{\rho_0}{r}$ is enclosed in a hollow sphere of radius $R_2$ with negative surface charge density $\sigma$, such that the total charge in the system is zero $\rho_0$ is a positive constant and $r$ is the distance from the centre of the sphere. The ratio $\dfrac{R_2}{R_1}$ is :-
\begin{qoptions}
\opt{A}{$\dfrac{\sigma}{\rho_0}$}
\opt{B}{$\sqrt{\dfrac{2\sigma}{\rho_0}}$}
\opt{C}{$\sqrt{\dfrac{\rho_0}{2\sigma}}$}
\opt{D}{$\dfrac{\rho_0}{\sigma}$}
\end{qoptions}
\qn{11} A uniform rod of mass $m=2$ kg and length $\ell=0.5$ m is sliding along two mutually perpendicular smooth walls with the two ends $P$ and $Q$ having velocities $v_P=4$ m/s and $v_Q=3$ m/s as shown. Then, choose the correct statement at the given instant ?
\qfig{p1_phy_s1_q11_1}
\begin{qoptions}
\opt{A}{The angular velocity of rod, $\omega=10$ rad/s. clockwise.}
\opt{B}{The angular velocity of rod $\omega=7.5$ rad/s. clockwise.}
\opt{C}{The velocity of centre of mass of rod. $v_{cm}=5$ m/s.}
\opt{D}{The total kinetic energy of rod, $K=\dfrac{25}{3}$ joule.}
\end{qoptions}
\qn{12} Charge $Q_0$ and $2Q_0$ are given to parallel plate A and B respectively and they are separated by a small distance. The capacitance of the given arrangement is $C$. Now plates A and B are connected to positive and negative terminals of battery of potential difference $V=\dfrac{2Q_0}{C}$ respectively as shown, then the work done by the battery is :-
\qfig{p1_phy_s1_q12_1}
\begin{qoptions}
\opt{A}{$\dfrac{2Q_0^2}{C}$}
\opt{B}{$\dfrac{4Q_0^2}{C}$}
\opt{C}{$\dfrac{5Q_0^2}{C}$}
\opt{D}{$\dfrac{6Q_0^2}{C}$}
\end{qoptions}
\qn{13} Four equal capacitors, each with capacitance $C$ are connected as shown in figure. Then the equivalent capacitance between A and B is
\qfig{p1_phy_s1_q13_1}
\begin{qoptions}
\opt{A}{$C$}
\opt{B}{$3C$}
\opt{C}{$2C$}
\opt{D}{$\dfrac{3C}{2}$}
\end{qoptions}
\qn{14} Two uniform boards, tied together with the help of string are balanced on a surface as shown. The coefficient of static friction between boards and surface is $0.5$. The minimum value for $\theta$, for which this type of arrangement is possible is :
\qfig{p1_phy_s1_q14_1}
\begin{qoptions}
\opt{A}{$30^\circ$}
\opt{B}{$45^\circ$}
\opt{C}{$37^\circ$}
\opt{D}{$53^\circ$}
\end{qoptions}
\qn{15} Two circular coils X and Y have equal number of turns and carry equal currents in the same sense and subtend same solid angle at point O. If the smaller coil X is midway between O and Y, then if we represent the magnetic induction due to bigger coil Y at O as $B_y$ and that due to smaller coil X at O as $B_x$, then:
\qfig{p1_phy_s1_q15_1}
\begin{qoptions}
\opt{A}{$\dfrac{B_y}{B_x}=1$}
\opt{B}{$\dfrac{B_y}{B_x}=2$}
\opt{C}{$\dfrac{B_y}{B_x}=\dfrac12$}
\opt{D}{$\dfrac{B_y}{B_x}=\dfrac14$}
\end{qoptions}
\qn{16} Gravitational field at a distance $2R$ from the centre of a spherical shell of inner radius $R$, outer radius $3R$ and density $d$ is equal to $k(G.\pi.R.d)$, where $G$ is universal gravitational constant and $k$ is
\begin{qoptions}
\opt{A}{$\dfrac73$}
\opt{B}{$\dfrac53$}
\opt{C}{$3$}
\opt{D}{$\dfrac43$}
\end{qoptions}
\qn{17} A thick walled hollow sphere has outer radius $R$. It rolls down an incline without slipping and its speed at bottom is $v_0$. Now the incline is waxed so that the friction is absent. The sphere is observed to slide down without rolling and the speed is now $\left(\dfrac{5v_0}{4}\right)$. The radius of gyration of the hollow sphere about the axis through its centre is
\begin{qoptions}
\opt{A}{$\dfrac34R$}
\opt{B}{$\dfrac{\sqrt3}{4}R$}
\opt{C}{$\dfrac23R$}
\opt{D}{None of these}
\end{qoptions}
\qn{18} An electric dipole has a fixed dipole moment $\vec p$, which makes angle $\theta$ with respect to $x$-axis. When subjected to an electric field $\vec E_1=E\hat i$, it experiences a torque $\vec T_1=\tau\hat k$. When subjected to another electric field $\vec E_2=\sqrt3E_1\hat j$ it experiences torque $\vec T_2=-\vec T_1$. The angle $\theta$ is :
\begin{qoptions}
\opt{A}{$60^\circ$}
\opt{B}{$90^\circ$}
\opt{C}{$30^\circ$}
\opt{D}{$45^\circ$}
\end{qoptions}
\qn{19} A cell develops the same power across two resistance $R_1$ and $R_2$ separately. The internal resistance of the cell is:
\begin{qoptions}
\opt{A}{$R_1+R_2$}
\opt{B}{$\dfrac{R_1+R_2}{2}$}
\opt{C}{$\sqrt{R_1R_2}$}
\opt{D}{$\dfrac{\sqrt{R_1+R_2}}{2}$}
\end{qoptions}
\qn{20} A thin uniform annular disc (sec figure) of mass $M$ has outer radius $4R$ and inner radius $3R$. The work required to take a unit mass from point $P$ on its axis to infinity is
\qfig{p1_phy_s1_q20_1}
\begin{qoptions}
\opt{A}{$\dfrac{2GM}{7R}\left(4\sqrt2-5\right)$}
\opt{B}{$-\dfrac{2GM}{7R}\left(4\sqrt2-5\right)$}
\opt{C}{$\dfrac{GM}{4R}$}
\opt{D}{$\dfrac{2GM}{5R}\left(\sqrt2-1\right)$}
\end{qoptions}
\sect{SECTION-II}
\qn{1} If a planet were suddenly stopped in its orbit supposed to be circular, it would fall into the sun in a time $T\left(\dfrac{\sqrt2}{n}\right)$, where $T$ is the time period of revolution. Find $6n$.
\qn{2} An infinite collection of current carrying infinite conductors each carrying a current $I$ outwards perpendicular to papers are placed at $x=\alpha_0,-2\alpha_0,3\alpha_0,-4\alpha_0\ldots.$at infinite here $\alpha_0$ is a positive constant. Then the magnetic field at the origin due to above collection of current carrying conductor is $\dfrac{\mu_0I\,\ell n8}{n\pi\alpha_0}$. Find $10n$.
\qn{3} A battery of 10 volt is connected to a resistance of 20 ohm through a variable resistance $R$, the amount of charge (in coulomb) which has passed in the circuit in 4 minutes, if the variable resistance $R$ is increased from zero at the rate of 5 ohm/min is $P\ell n2$. Find $P$.
\qn{4} In the circuit shown, the ratio of power in $8\Omega$ to that in $4\Omega$ is $n$. Find $n$.
\qfig{p1_phy_s2_q04_1}
\qn{5} A uniform thin rod AB is equipped at both ends with the hooks shown and is supported by a frictionless horizontal table. Initially the rod is hooked at A to a fixed pin C about which it rotates with a constant angular velocity $\omega_1$. Suddenly end B of the rod hits pin D and gels hooked to pin D. causing end A to be released. The magnitude of the angular velocity $\omega_2$ of the rod in its subsequent rotation about D is $\dfrac{\omega_1}{P}$. Find $3P$. (Assume length and mass of the hook is negligible. Pin C \& D are lying on a same horizontal line)
\qfig{p1_phy_s2_q05_1}
\qn{6} A very long cylinder is carrying current $I_0$ distributed uniformly over its cross-section area, as shown in the figure. the value of $\displaystyle\int\vec B.d\vec\ell$ along the path AB (from A to B) shown in the figure. $O$ is the centre of the cross section of the cylinder and direction of current is in to the plane of the paper is $\dfrac{-\mu.i.n}{78}$.
\qfig{p1_phy_s2_q06_1}
\qn{7} In the circuit shown, the charge ($\mu$C) in the 3 uF capacitor at steady state will be :-
\qfig{p1_phy_s2_q07_1}
\qn{8} A cell of emf $V$ is connected to a capacitor having square plates of dimension $a$ and plate separation $d$. A dielectric slab of dielectric constant $(\sqrt3+1)$ and thickness $d$ is placed between the plates such that it exactly fills up the space between the plates at $t=0$. Now the slab is rotated with angular velocity $\omega$ as shown in the figure. Neglect edge effects the capacitance at $t=\dfrac{\pi}{6\omega}$ is $\dfrac{x\varepsilon.a^2}{8d}$. Find $x$.
\qfig{p1_phy_s2_q08_1}
\qn{9} A slab of resistance $R$ is inserted between two parallel plates of a capacitor charged to $Q_0$. The capacitor is discharged through the solid. A magnetic field is present throughout and is uniform and out of the plane of paper. The total momentum given to the slab after complete discharge is $nBQ.d$. Find $n$ (Given capacitance is $C$ and plates are separated by $d$).
\qfig{p1_phy_s2_q09_1}
\qn{10} In a satellite the time of revolution is $T$. Then K.E. is proportional to $T^n$. Find $(-24\ n)$
\qfullbegin
\partsec{2}{CHEMISTRY}
\sect{SECTION-I}
\qn{1} Select the correct graph between magnetic moment of complex compound and number of d-electrons with weak field ligand and strong field ligands respectively.
\begin{qoptions}
\opt{A}{\qfigi{p1_chem_s1_q01_A}}
\opt{B}{\qfigi{p1_chem_s1_q01_B}}
\opt{C}{\qfigi{p1_chem_s1_q01_C}}
\opt{D}{\qfigi{p1_chem_s1_q01_D}}
\end{qoptions}
\qfullend
\qn{2} Which of the following is not true for metal carbonyls?
\begin{qoptions}
\opt{A}{The oxidation state of central metal in the carbonyls can be zero, negative or low positive.}
\opt{B}{d$\pi$--d$\pi$ back bonding from M $\to$ CO is present}
\opt{C}{d$\pi-\pi*$ MO back bonding is present from M $\to$ CO}
\opt{D}{bond order of CO decreases during M $\xrightarrow{\pi}$ CO}
\end{qoptions}
\qn{3} Total number of isomers of a compound having molecular formula, $\mathrm{[Pd(en)_2(NH_3)(H_2O)]^{4+}}$ is :
\begin{qoptions}
\opt{A}{$0$}
\opt{B}{$1$}
\opt{C}{$2$}
\opt{D}{$3$}
\end{qoptions}
\qn{4} In which of the following complex, number of stereoisomers are less than the number the structural isomers by 2?
\begin{qoptions}
\opt{A}{$\mathrm{K_3[Fe(C_2O_4)_3]}$}
\opt{B}{$\mathrm{[Pt(gly)_2]}$}
\opt{C}{$\mathrm{[Co(NH_3)_5(NO_2)](NO_3)_2}$}
\opt{D}{$\mathrm{[Co(NH_3)_4(NO_2)_2]^+}$}
\end{qoptions}
\qn{5} Select the incorrect statement :
\begin{qoptions}
\opt{A}{Addition elimination mechanism is involved in the hydrolysis of both $\mathrm{SO_2Cl_2}$ and $\mathrm{SOCl_2}$}
\opt{B}{The bond length decreases from $\mathrm{NO\to NO^+}$ and also decreases from $\mathrm{CO\to CO^+}$}
\opt{C}{In $\mathrm{PCl_2F_3}$, number of bond angle(s) less than $120^\circ$ but greater than $90^\circ$ are 2}
\opt{D}{The bond angle of $\mathrm{OCl_2}$ is greater than dimethyl ether due to presence of back bonding}
\end{qoptions}
\qn{6} Which of the following is incorrect
\begin{qoptions}
\opt{A}{$\mathrm{H_2>H}$ (ionization energy)}
\opt{B}{$\mathrm{O_2<O}$ (ionization energy)}
\opt{C}{$\mathrm{He<He^+}$ (ionization energy)}
\opt{D}{$\mathrm{N_2<N_2^+}$ (bond order)}
\end{qoptions}
\qn{7} The kinetics of an enzyme-catalyzed hydrolysis reaction is represented by the following graph, where $[S]$ is the substrate concentration, and $v$ is the rate of reaction :
\qfig{p1_chem_s1_q07_1}
The kinetic course of the reaction can be described as :
\begin{qoptions}
\opt{A}{First order, zero order.}
\opt{B}{Zero order, first order}
\opt{C}{First order, second order}
\opt{D}{Second order, first order}
\end{qoptions}
\qn{8} FeO crystallises as NaCl structure. Due to same metal deficient defect its formula is found as $\mathrm{Fe_{0.75}O}$. The percent decrease in density of solid due to defect is : (M. wt. of Fe $=56$)
\begin{qoptions}
\opt{A}{$25\%$}
\opt{B}{$19.44\%$}
\opt{C}{$75\%$}
\opt{D}{$12.5\%$}
\end{qoptions}
\qn{9} Decomposition of a non-volatile solute A into another non-volatile solute B and C in an aqueous solution follows first-order kinetics as : $\mathrm{A\to2B+C}$\\
When one mole of A is dissolved in 180 gm water and left for decomposition, the vapour pressure of solution was found to be 20 mm Hg after 12 hr. What is the vapour pressure of solution after 24 hr. Assume $25^\circ$C temperature throughout and the vapour pressure of water at $25^\circ$C is 24 mm Hg.
\begin{qoptions}
\opt{A}{$21.6$ mm Hg}
\opt{B}{$19.2$ mm Hg}
\opt{C}{$10.8$ mm Hg}
\opt{D}{$16$ mm Hg}
\end{qoptions}
\qn{10} If an ionic solid of cation A and anion B has radius ratio of cation to anion equal to $\left(\sqrt{\dfrac52}-1\right)$, then how many anions does each cation touch :
\begin{qoptions}
\opt{A}{$4$}
\opt{B}{$6$}
\opt{C}{$12$}
\opt{D}{$3$}
\end{qoptions}
\qn{11} Two solutions $S_1$ and $S_2$ containing 0.1 M NaCl (aq) and 0.05 M $\mathrm{BaCl_2}$ (aq) are sparated by semiperameable membrane. Which among the following statement(s) is/are correct ? (Assume complete dissociation of both the electrolytes) :-
\qfig{p1_chem_s1_q11_1}
\begin{qoptions}
\opt{A}{$S_1$ and $S_2$ are isotonic}
\opt{B}{$S_1$ is hypertonic while $S_2$ is hypotonic}
\opt{C}{$S_1$ is hypotonic while $S_2$ is hypertonic}
\opt{D}{All the above}
\end{qoptions}
\qn{12} The Henry's law constant for the two gases A and B are $6.4\times10^4$ atm and $4\times10^4$ atm respectively in water. If ratio of mole fractions of A to B in gaseous mixture above the solution is 0.2, then find mole ratio of of B to A dissolved in the solution.
\begin{qoptions}
\opt{A}{$8$}
\opt{B}{$6.4$}
\opt{C}{$1.6$}
\opt{D}{$0.4$}
\end{qoptions}
\qn{13} Select the order which is not correct among following :
\begin{qoptions}
\opt{A}{\qfigi{p1_chem_s1_q13_A}}
\opt{B}{\qfigi{p1_chem_s1_q13_B}}
\opt{C}{\qfigi{p1_chem_s1_q13_C}}
\opt{D}{\qfigi{p1_chem_s1_q13_D}}
\end{qoptions}
\qn{14} Choose the halogen which is most reactive towards $\mathrm{S_N1}$ reaction in the given compounds (A, B, C and D)
\begin{qoptions}
\opt{A}{\qfigi{p1_chem_s1_q14_A}}
\opt{B}{\qfigi{p1_chem_s1_q14_B}}
\opt{C}{\qfigi{p1_chem_s1_q14_C}}
\opt{D}{\qfigi{p1_chem_s1_q14_D}}
\end{qoptions}
\begin{qoptions}
\opt{A}{A--Br$_{\rm(b)}$, B--I$_{\rm(b)}$, C--Br$_{\rm(b)}$, D--Br$_{\rm(b)}$}
\opt{B}{A--Br$_{\rm(a)}$, B--I$_{\rm(b)}$, C--Br$_{\rm(b)}$, D--Br$_{\rm(a)}$}
\opt{C}{A--Br$_{\rm(b)}$, B--I$_{\rm(b)}$, C--Br$_{\rm(a)}$, D--Br$_{\rm(a)}$}
\opt{D}{A--Br$_{\rm(b)}$, B--I$_{\rm(a)}$, C--Br$_{\rm(a)}$, D--Br$_{\rm(a)}$}
\end{qoptions}
\qn{15} In the following reaction, B is :
\qfig{p1_chem_s1_q15_1}
\begin{qoptions}
\opt{A}{\qfigi{p1_chem_s1_q15_A}}
\opt{B}{\qfigi{p1_chem_s1_q15_B}}
\opt{C}{\qfigi{p1_chem_s1_q15_C}}
\opt{D}{\qfigi{p1_chem_s1_q15_D}}
\end{qoptions}
\qn{16} Consider the following reaction that goes from A to B in three steps as shown below :
\qfig{p1_chem_s1_q16_1}
Choose the correct\\
\begin{center}
{\small\setlength{\tabcolsep}{3pt}\begin{tabular}{|p{0.22\linewidth}|p{0.25\linewidth}|p{0.22\linewidth}|}\hline
Number of intermediates & Number of Activated complexes & Rate determining step\\\hline
\end{tabular}}
\end{center}
\begin{qoptions}
\opt{A}{$3$ \quad $2$ \quad II}
\opt{B}{$2$ \quad $3$ \quad II}
\opt{C}{$2$ \quad $3$ \quad I}
\opt{D}{$2$ \quad $3$ \quad III}
\end{qoptions}
\qn{17} Match List I with List II\\
I-Bromopropane is reacted with reagents in List I to give product in List II
\begin{center}
\begin{tabular}{|c|l|c|l|}\hline
 & \textbf{List-I Reagent} & & \textbf{List-II Product}\\\hline
A & KOH (alc) & I & Nitrile\\\hline
B & KCN (alc) & II & Ester\\\hline
C & $\mathrm{AgNO_2}$ & III & Alkene\\\hline
D & $\mathrm{H_3CCOOAg}$ & IV & Nitroalkane\\\hline
\end{tabular}
\end{center}
\begin{qoptions}
\opt{A}{A-IV, B-III, C-II, D-I}
\opt{B}{A-III, B-I, C-IV, D-II}
\opt{C}{A-I, B-II, C-III, D-IV}
\opt{D}{A-I, B-III, C-IV, D-II}
\end{qoptions}
\qn{18} 2-Methyl propyl bromide reacts with $\mathrm{C_2H_6O^-}$ and gives A whereas on reaction with $\mathrm{C_2H_5OH}$ it gives B. The mechanism followed in these reactions and the products A and B respectively are :
\begin{qoptions}
\opt{A}{$\mathrm{S_N2}$, A = iso-butyl ethyl ether, $\mathrm{S_N1}$, B = tert-butyl ethyl ether}
\opt{B}{$\mathrm{S_N1}$, A = tert-butyl ethyl ether, $\mathrm{S_N1}$, B = 2-butyl ethyl ether}
\opt{C}{$\mathrm{S_N1}$, A = tert-butyl ethyl ether, $\mathrm{S_N2}$, B = iso-butyl ethyl ether}
\opt{D}{$\mathrm{S_N2}$, A = 2-butyl ethyl ether, $\mathrm{S_N2}$, B = iso-butyl ethyl ether}
\end{qoptions}
\qn{19} Which of the following Fischer projection is the enantiomer of the following molecule?
\qfig{p1_chem_s1_q19_1}
\begin{qoptions}
\opt{A}{\qfigi{p1_chem_s1_q19_A}}
\opt{B}{\qfigi{p1_chem_s1_q19_B}}
\opt{C}{\qfigi{p1_chem_s1_q19_C}}
\opt{D}{\qfigi{p1_chem_s1_q19_D}}
\end{qoptions}
\qn{20} {\footnotesize\setlength{\tabcolsep}{2pt}\begin{tabular}[t]{|c|p{0.3\linewidth}|c|p{0.28\linewidth}|}\hline
 & \textbf{List-I (Complexes)} & & \textbf{List-II (Properties)}\\\hline
(a) & $\mathrm{[Co(H_2O)_6]^{3+}}$ & (i) & $-2.4\Delta_0$\\\hline
(b) & $\mathrm{[Mn(NH_3)_6]^{2+}}$ & (ii) & Contains unsymmetric bidented ligand\\\hline
(c) & $\mathrm{[Zn(gly)_2]}$ & (iii) & paramagnetic\\\hline
(d) & $\mathrm{[IrF_6]^{3-}}$ & (iv) & $d^2sp^3$ hybridisation\\\hline
\end{tabular}}\par
Select \textbf{CORRECT} matching :-
\begin{qoptions}
\opt{A}{a--(iii) ; b--(iv) ; c--(iv) ; d--(i)}
\opt{B}{a--(iv) ; b--(i) ; c--(iii) ; d--(iii)}
\opt{C}{a--(iv) ; b--(iii) ; c--(ii) ; d--(i)}
\opt{D}{a--(iv) ; b--(i) ; c--(ii) ; d--(iii)}
\end{qoptions}
\sect{SECTION-II}
\qn{1} Number of molecules which are polar and planar.\\
$\mathrm{BF_3}$, $\mathrm{B_2H_6}$, $\mathrm{I_2Cl_6}$, $\mathrm{SO_3}$, $\mathrm{SO_2F_2}$, $\mathrm{XeO_3}$, $\mathrm{ClO_3^-}$, $\mathrm{XeOF_4}$
\qn{2} Select the ligands which form more than 2 rings upon acting as ligand.\\
Dien, o-phen, pn, trien, $\mathrm{NTA^{3-}}$, $\mathrm{acac^-}$, $\mathrm{Imda^{2-}}$, $\mathrm{EDTA^{4-}}$
\qn{3} Compounds:
\qfig{p1_chem_s2_q03_1}
(ix) $\mathrm{CH_3{-}CH_2{-}OH}$\\
$x$ = number of compounds soluble in NaOH (aq.)\\
$y$ = number of compounds soluble in $\mathrm{NaHCO_3}$ (aq.)\\
then $x+y=$ ?
\qn{4} $X$ = number of following alcohols give immidiate turbidity with Lucas reagent.\\
Find value of $X$
\qfig{p1_chem_s2_q04_1}
(iii) $\mathrm{CH_3{-}OH}$ \quad (iv) $\mathrm{CH_3{-}CH_2{-}OH}$
\qfig{p1_chem_s2_q04_2}
\qfig{p1_chem_s2_q04_3}
\qn{5} The graph below represents the distribution of molecular speeds of gas A and gas B at 200 K. If equimolar mixture of gases A and B at 200 K is effuses through small hole then calculate the relative rate of effusion of gas B relative to gas A at 200 K.
\qfig{p1_chem_s2_q05_1}
\qn{6} 1.6 gram of weak acid $\mathrm{CH_2FCOOH}$ is dissolved in water preparing 500 ml aqueous solution. Find osmotic pressure of this solution in atm if the acid is 60\% ionized in water at 300 K ?\\
{[}Given $R=0.082$ L-atm mol$^{-1}$ K$^{-1}${]}\\
{[}Atomic mass : F $=19${]}
\qn{7} For an ideal gas sample (fixed mole taken in fixed volume container) collision frequency ($Z$) \& mean free path ($\lambda$).\\
(i) Both $Z$ \& $\lambda$ change with change in $T$.\\
(ii) $Z$ changes in $T$.\\
(iii) $\lambda$ changes in $T$.\\
(iv) Both are independent of temperature $T$.\\
(v) Heavier gas (more mole weight) exibits high $Z$ compared to lighter gas, at some temperature.\\
(vi) Heavier gas exibits more $\lambda$ compared to lighter gas at same temperature.\\
(vii) Molecular weight of gas does not have any effect of value of $\lambda$.\\
Identify number of correct statements from above ?
\qn{8} Reaction $\mathrm{A+B\to C+D}$ follow's following rate law : rate $=k=[A]^{+\frac12}[B]^{\frac12}$. Starting with initial conc. of one mole of A and B each, what is the time taken for amount of A to become 0.25 mole. [Given $k=2.31\times10^{-3}$ sec$^{-1}$]\\
($\ln2=0.693$)
\qn{9} Number of meso isomer of the given compound is :
\qfig{p1_chem_s2_q09_1}
\qn{10} How many of the following elements can be purified by Van Arkel method.\\
Ti, Zr, Hf, Fe, Pb, Sn, B
\partsec{3}{MATHEMATICS}
\sect{SECTION-I}
\qn{1} If $f(x)=\displaystyle\lim_{h\to0}\dfrac{e^{\sin^{-1}(x+h)}-e^{\sin^{-1}x}}{h}$, then equation of tangent at $x=0$ to the curve is given by -
\begin{qoptions}
\opt{A}{$y=2x+1$}
\opt{B}{$y=x+3$}
\opt{C}{$2x+y=1$}
\opt{D}{$y=x+1$}
\end{qoptions}
\qn{2} If $y^2=2ax$, $a>0$ and $xy=4\sqrt2$ intersect orthogonally, then $a$ equals to
\begin{qoptions}
\opt{A}{$\dfrac12$}
\opt{B}{$2$}
\opt{C}{$\dfrac32$}
\opt{D}{$4$}
\end{qoptions}
\qn{3} Let $f(x)=xe^{-x(x+1)}$ is increasing in $\left(-a,\dfrac ab\right)$, where $a\in N$, then value of $(a+b)$ is
\begin{qoptions}
\opt{A}{$3$}
\opt{B}{$1$}
\opt{C}{$\dfrac12$}
\opt{D}{$2$}
\end{qoptions}
\qn{4} Range of $f(x)=k\sin x+\cos x$ is $[-5,5]$, then total number of integers between the two possible value of $k$ is -
\begin{qoptions}
\opt{A}{$11$}
\opt{B}{$10$}
\opt{C}{$9$}
\opt{D}{$5$}
\end{qoptions}
\qn{5} If $\sin A+\cos B=a$ and $\sin B+\cos A=b$, then $\sin(A+B)$ is equal to -
\begin{qoptions}
\opt{A}{$\dfrac{a^2+b^2}{2}$}
\opt{B}{$\dfrac{a^2-b^2+2}{2}$}
\opt{C}{$\dfrac{a^2+b^2-2}{2}$}
\opt{D}{$\dfrac{a^2+b^2+2}{2}$}
\end{qoptions}
\qn{6} The value of $\cos\dfrac\pi5\cos\dfrac{2\pi}5\cos\dfrac{4\pi}5\cos\dfrac{8\pi}5$ is equal to -
\begin{qoptions}
\opt{A}{$\dfrac1{16}$}
\opt{B}{$\dfrac14$}
\opt{C}{$-\dfrac18$}
\opt{D}{$-\dfrac1{16}$}
\end{qoptions}
\qn{7} Number of solutions of the equation $\tan\theta+\tan2\theta+\tan3\theta=\tan\theta\tan2\theta\tan3\theta$ in $(0,\pi)$ is -
\begin{qoptions}
\opt{A}{$5$}
\opt{B}{$3$}
\opt{C}{$2$}
\opt{D}{zero}
\end{qoptions}
\qn{8} The number of solutions of $2^{\sin|x|}=4^{|\cos x|}$ in $[-\pi,\pi]$ is -
\begin{qoptions}
\opt{A}{$0$}
\opt{B}{$2$}
\opt{C}{$4$}
\opt{D}{$6$}
\end{qoptions}
\qn{9} The general solution of the equation $\sin x-\cos x=\min\left\{\dfrac1{\sqrt2},x^2-6x+10\right\}$ where $n\in I$, is -
\begin{qoptions}
\opt{A}{$\dfrac\pi4+n\pi+(-1)^n\dfrac\pi6$}
\opt{B}{$n\pi+(-1)^n\dfrac\pi6-\dfrac\pi4$}
\opt{C}{$2n\pi$}
\opt{D}{$2n\pi+\dfrac\pi2$}
\end{qoptions}
\qn{10} As shown in the figure, the area of the trapezium ABCD is $\dfrac{81}{2}$, where $a=2b$, then $c$ is equal to -
\qfig{p1_math_s1_q10_1}
\begin{qoptions}
\opt{A}{$\dfrac92$}
\opt{B}{$3\sqrt5$}
\opt{C}{$\sqrt{30}$}
\opt{D}{none of these}
\end{qoptions}
\qn{11} A value of $c$ for which the Rolle's theorem holds for the function $f(x)=2x^3+x^2-4x-2$ on the interval $[-\sqrt2,\sqrt2]$ is
\begin{qoptions}
\opt{A}{$\dfrac35$}
\opt{B}{$\dfrac45$}
\opt{C}{$\dfrac23$}
\opt{D}{$-\dfrac23$}
\end{qoptions}
\qn{12} Suppose that $f$ is differentiable function such that $f'(x)\leq2$ for all $x$. If $f(1)=2$ and $f(4)=8$ then $f(2)$ has the value equal to
\begin{qoptions}
\opt{A}{$4$}
\opt{B}{$2$}
\opt{C}{$1$}
\opt{D}{$0$}
\end{qoptions}
\qn{13} A spherical balloon is filled with $7500\pi$ m$^3$ of helium gas. If a leakage in the balloon starts at the rate of $100\pi$ m$^3$/min., then the rate (in m/min) at which the radius of balloon will be decreasing after 30 minutes
\begin{qoptions}
\opt{A}{$\dfrac49$}
\opt{B}{$\dfrac29$}
\opt{C}{$\dfrac19$}
\opt{D}{$\dfrac13$}
\end{qoptions}
\qn{14} Which of the following inequality is incorrect ?
\begin{qoptions}
\opt{A}{$\ell n\left(1+\dfrac1x\right)>\dfrac{2}{2x+1}$ for $x>0$}
\opt{B}{$\sin x<x<\tan x$ for $x\in\left(0,\dfrac\pi2\right)$}
\opt{C}{$2x\sec x+x>3\tan x$ for $x\in\left(0,\dfrac\pi2\right)$}
\opt{D}{$\sin x<\dfrac{2x}{\pi}$ for $x\in\left(0,\dfrac\pi2\right)$}
\end{qoptions}
\qn{15} The function $f(x)=\begin{cases}\dfrac{|x|-x\left(3^{1/x}+1\right)}{3^{1/x}-1} & x\neq0\\ 0 & x=0\end{cases}$ is
\begin{qoptions}
\opt{A}{discontinuous at $x=0$}
\opt{B}{continuous at $x=0$ but not differentiable at $x=0$}
\opt{C}{both continuous and differentiable at $x=0$}
\opt{D}{differentiable but not continuous at $x=0$}
\end{qoptions}
\qn{16} The value of $\displaystyle\lim_{\theta\to0^+}\left(\theta^2\sqrt\theta\cdot\sin\left(\dfrac1\theta\right)+4\sin(\theta\sqrt\theta)\right)^\theta$ is
\begin{qoptions}
\opt{A}{$-1$}
\opt{B}{$1$}
\opt{C}{$2$}
\opt{D}{$e$}
\end{qoptions}
\qn{17} $f(x)=x^5+e^{x/3}$, $g(x)=f^{-1}(x)$ then $g'(1)=$
\begin{qoptions}
\opt{A}{$2$}
\opt{B}{$1$}
\opt{C}{$3$}
\opt{D}{None of these}
\end{qoptions}
\qn{18} $y=e^{2\sqrt x}+e^{-2\sqrt x}$, then $\left(x\dfrac{d^2y}{dx^2}+\dfrac12\dfrac{dy}{dx}\right)+2$ equals to
\begin{qoptions}
\opt{A}{$y+2$}
\opt{B}{$xy+2$}
\opt{C}{$\dfrac{y+8}{4}$}
\opt{D}{$\sqrt{xy}$}
\end{qoptions}
\qn{19} Let
\[f(x)=\begin{cases}x^2e^{2(x-1)}; & 0\leqslant x\leqslant1\\ a\,\mathrm{sgn}\,(x+1)\cos(2x-2)+bx^2, & 1<x\leqslant2\end{cases}.\]
If $f(x)$ is differentiable at $x=1$, then (where $\mathrm{sgn}(x)$ denotes signum function of $x$)
\begin{qoptions}
\opt{A}{$a=-1$, $b=2$}
\opt{B}{$a=-3$, $b=-2$}
\opt{C}{$a=-3$, $b=4$}
\opt{D}{$a=-5$, $b=-4$}
\end{qoptions}
\qn{20} Which of the following function is discontinuous at $x=0$, if $f(0)=1$
\begin{qoptions}
\opt{A}{$f(x)=\dfrac{\sin(\sin^2x)}{\left(e^{|\sin x|}-1\right)\sqrt{x^2}}$; $(x\neq0)$}
\opt{B}{$f(x)=(1+x)(\mathrm{sgn}(x))+x^2$; $(x\neq0)$ (where $\mathrm{sgn}(x)$ is signum function)}
\opt{C}{$f(x)=(\ell n(1+\tan^2x))(\mathrm{cosec}((e^x-1)^2))$; $(x\neq0)$}
\opt{D}{$f(x)=(1+\sin^2x)^{\mathrm{cosec}\,x}$; $(x\neq0)$}
\end{qoptions}
\sect{SECTION-II}
\qn{1} If $f(x)=10x-2x^2$, then {\small$\displaystyle\lim_{h\to0}\dfrac{f(3\cosh+4\sinh-2)-8}{f(3e^h-5\sec h+4)-12}$} is equal to
\qn{2} If $f(x)=\begin{cases}\dfrac{e^{ax+\ell n(1+4x)+b}}{x^3} & x\neq0\\ L & x=0\end{cases}$ is continuous at $x=0$, then $12abL$ is equal to
\qn{3} Let $f(x)=(x-1)^{100}(x-2)^{2(99)}(x-3)^{3(98)}\ldots.(x-100)^{100}$. If $P=\dfrac{f'(101)}{f(101)}$, then $\dfrac P{50}$ equals
\qn{4} Number of integers in the range of function $f(x)=\dfrac{100}{3(x^6+12x^4+5x^2+1)}$ is
\qn{5} If the solution set of inequality $\left(\cos x+\dfrac{\sqrt3}{2}\right)\left(\cos x-\dfrac12\right)\leqslant0$ in $[0,2\pi]$ is $\left[\dfrac{\alpha\pi}6,\dfrac{\beta\pi}6\right]\cup\left[\dfrac{\gamma\pi}6,\dfrac{\delta\pi}6\right]$, $\alpha,\beta,\gamma,\delta\in I^+$, then the value of $|\beta-\alpha+\delta-\gamma|$ is
\qn{6} If $x\in\left(\pi,\dfrac{3\pi}2\right)$, then value of $4\cos^2\left(\dfrac\pi4-\dfrac x2\right)+\sqrt{4\sin^4x+\sin^22x}$ is equal to
\qn{7} If $A+B+C=\pi$, and $\dfrac{\Sigma\sin2A}{\Sigma\sin A}=\lambda\displaystyle\prod\sin\dfrac A2$, then value of $\lambda$ is
\qn{8} If $f(x)+f\left(1-\dfrac1x\right)=1+x\ \forall\ x\in R-\{0,1\}$, then value of $16f(2)$ is
\qn{9} For $0<x<1$, if the solution of equation $3\sin^{-1}\left(\dfrac{2x}{1+x^2}\right)+2\tan^{-1}\left(\dfrac{2x}{1-x^2}\right)-4\cos^{-1}\left(\dfrac{1-x^2}{1+x^2}\right)=\dfrac\pi4$ is $K$ then value of $100K$ is (use $\sqrt2=1.41$)
\qn{10} If $\displaystyle\lim_{x\to2}\dfrac{1-\cos\left(x^2-6x+8\right)}{(x-2)^2}=p$, then value of $\left(\dfrac6p\right)$ is equal to
\papersec{2}
\partsec{1}{PHYSICS}
\sect{SECTION-I(i)}
\qn{1} A circuit diagram of an electronics device's part is given in the figure. The diagram consists of a 9 V source, a filter made of a capacitor and a resistor, a switch and output leads. The two possible positions of the switch are "on (A and B are connected) and "off" (A and C are connected). If in the given circuit the switch is turned from the position "on" to "off" (that means to turn the device off) and next the voltage source's polarity is reversed then the electronics device will still work after turning it on. However, if the voltage source's polarity is changed without turning the device off, then the resistor $R$ will fuse. Assuming that the resistor will fuse when the power on it exceeds $0.25$ W, choose the correct option(s) :-
\qfig{p2_phy_s1_q01_1}
\begin{qoptions}
\opt{A}{The minimum value of $R$ is $400\Omega$}
\opt{B}{The minimum value of $R$ is $324\Omega$}
\opt{C}{The maximum value of $R$ is $1296\Omega$}
\opt{D}{The minimum value of $R$ is $1600\Omega$}
\end{qoptions}
\qn{2} Electric field intensity at a point $(x,y)$ is given by $\vec I=\left(12xy^3-4x\right)\hat i+18x^2y^2\hat j$ where all parameters are in S.I. Units. Then choose the correct statement(s) out of the following
\begin{qoptions}
\opt{A}{it represents a conservative electric field.}
\opt{B}{it represents a non-conservative electric field.}
\opt{C}{electric potential at point $(1\mathrm m,1\mathrm m)$ is $-4$V is electric potential at origin is zero.}
\opt{D}{electric potential at point $(1\mathrm m,1\mathrm m)$ cannot be defined if potential is taken zero at origin.}
\end{qoptions}
\qn{3} A positively charged particle is accelerated from the origin with the help of a voltage $4U$ (where $U>0$) to a certain velocity which lies on the $x-y$ plane under the angle of 30 degrees with respect to the $x$-axis (see figure). The electric potential $V(x,y)=V(x)$ only depends on the $x$-coordinate: if $3a<x<4a$ then $V\{x\}=U$ and if not then $V(x)=0$; besides the electrostatic force no other forces are applied to the particle and $a>0$.\\
For option C \& D consider the source of the charged particles is now a coaxial vacuum diode at the origin. Because of that the particles accelerated with the voltage $4U$ are moving isotropically (at equal quantities) to all directions of the $x-y$ plane; the $z$-directional component of the velocity for all particles is 0.
\qfig{p2_phy_s1_q03_1}
\begin{qoptions}
\opt{A}{The trajectory of the charge on x-y plane will be \qfigi{p2_phy_s1_q03_2}.}
\opt{B}{The trajectory of the charge on x-y plane will be \qfigi{p2_phy_s1_q03_3}.}
\opt{C}{The part of all the emitted particles that reaches the region of space $x>4a$ is $\dfrac13$.}
\opt{D}{The part of all the emitted particles that reaches the region of space $x>4a$ is $\dfrac23$.}
\end{qoptions}
\qn{4} A galvanometer is modified so that it can measure a physical quantity in the circuit. In the modification instead of using permanent magnet, current coil is used on a ferromagnet and the arrangement is such that magnetic field due to current coil always remains perpendicular to the area of voltage coil. The remaining arrangement of the galvanometer remains same. The modified circuit is shown below. Choose the correct option(s) :
\qfig{p2_phy_s1_q04_1}
\begin{qoptions}
\opt{A}{If the marking of the scale is uniform, it will measure power dissipated in the resistor.}
\opt{B}{If the marking of the scale is uniform, it will measure voltage drop in the resistor.}
\opt{C}{For increasing the sensitivity of the device, the frame on which the current coil is wrapped should be of ferro-magnetic material.}
\opt{D}{On increasing the value of current to twice deflection on the scale will also be doubled.}
\end{qoptions}
\qn{5} The figure shows a square wire mesh made of conducting wires. Here, each small square has side $a$. The structure is kept in uniform magnetic field $B$ perpendicular to the plane of paper:
\qfig{p2_phy_s1_q05_1}
\begin{qoptions}
\opt{A}{The magnetic force on the structure is $2\sqrt2\,iBa$}
\opt{B}{The potential of point B $=$ potential of point D}
\opt{C}{Potential of point O $=$ potential of point B}
\opt{D}{The magnetic force on the structure is $\sqrt2\,iBa$}
\end{qoptions}
\qn{6} Five equal resistances, a capacitor and two batteries are connected in a circuit as shown in the figure. The capacitor is initially uncharged and each battery has negligible internal resistance. It is given that $R=1\Omega$, $C=1\ \mu$F, $C=1\ \mu$F and $V=8$ volt. On the basis of above information answer the following questions, if switch S is closed at $t=0$.
\qfig{p2_phy_s1_q06_1}
\begin{qoptions}
\opt{A}{Current in the branch of capacitor at $t=5/8\ \mu$s is (approximately) 6A.}
\opt{B}{Current in the branch AB at $t=5/8\ \mu$s is (approximately) 1.7A.}
\opt{C}{Potential difference between points A and C at $t=5/8\ \mu$s is (approximately) 2.5 V}
\opt{D}{Potential difference between points A and C at $t=5/8\ \mu$s is (approximately) 6V.}
\end{qoptions}
\sect{SECTION-I(ii)}
\qn{7} List-I describes some situations in which a small object moves. List-II describes some characteristics of these motions. Match the situations in List-I with the characteristics in List-II
\begin{center}
{\footnotesize\setlength{\tabcolsep}{2pt}\begin{tabular}{|c|p{0.36\linewidth}|c|p{0.28\linewidth}|}\hline
\multicolumn{2}{|c|}{List-I} & \multicolumn{2}{c|}{List-II}\\\hline
(P) & The object moves on the $x$--axis under a conservative force in such a way that its "speed" and "position" satisfy $v=c_1\sqrt{c_2-x^2}$, where $c_1$ and $c_2$ are positive constants. & (1) & The object executes a simple harmonic motion.\\\hline
(Q) & The object moves on the $x$-axis in such a way that its velocity and its displacement from the origin satisfy $v=-kx$, where $k$ is a positive constant. & (2) & The object does not change its direction,\\\hline
(R) & The object is attached to one end of a massless spring of a given spring constant. The other end of the spring is attached to the ceiling of an elevator. Initially everything is at rest. The elevator starts going upwards with a constant acceleration. The motion of the object is observed from the elevator during the period it maintains this acceleration. & (3) & The kinetic energy of the object keeps on decreasing.\\\hline
(S) & The object is projected from the earth's surface vertically upwards with a speed $2\sqrt{GM_e/R_e}$, where $M_e$ is the mass of the earth and $R_e$ is the radius of the earth. Neglect forces from objects other than the earth. & (4) & The object can change its direction only once.\\\hline
\end{tabular}}
\end{center}
\begin{qoptions}
\opt{A}{$\mathrm P\to1;Q\to3,4;R\to2;S\to2,3$}
\opt{B}{$\mathrm P\to2;Q\to2,4;R\to1;S\to2,3$}
\opt{C}{$\mathrm P\to1;Q\to2,3;R\to1;S\to2,3$}
\opt{D}{$\mathrm P\to1;Q\to2,3;R\to4;S\to1,2$}
\end{qoptions}
\qn{8} \textbf{Answer the following by appropriately matching the lists based on the information given in the paragraph.}\\
A capacitor has a capacitance of 4 $\mu$F. It is filled with air and has an inter plate distance at $d$. We have a dielectric slab ($k=2$) of same thickness as the distance between the plates of the capacitor but whose area is half that of the capacitor. We perform four actions on the capacitor (i) Fill the capacitor with the dielectric, (ii) remove the dielectric slab (iii) Connect the capacitor to a battery of emf 4 V. (iv) Disconnect the battery. In column I, we mention the order of the actions. In column II, we list the various parameters after the actions. Match the entries.
\begin{center}
{\footnotesize\setlength{\tabcolsep}{2pt}\begin{tabular}{|c|c|c|p{0.4\linewidth}|}\hline
\multicolumn{2}{|c|}{List-I} & \multicolumn{2}{c|}{List-II}\\\hline
(P) & (i)(iii)(ii)(iv) & (1) & Final Charge on capacitor $=16\mu$C\\\hline
(Q) & (iii)(iv)(i)(ii) & (2) & Final energy stored in the capacitor $=32\mu$J\\\hline
(R) & (i)(iii)(iv)(ii) & (3) & Final voltage across the capacitor $=6$V\\\hline
(S) & (iii)(i)(iv)(ii) & (4) & Final Charge on capacitor $=24\ \mu$C\\\hline
 & & (5) & Final energy stored in the capacitor $=72\mu$J\\\hline
 & & (6) & Final voltage across the capacitor is 4V.\\\hline
\end{tabular}}
\end{center}
\begin{qoptions}
\opt{A}{$\mathrm P\to1,5,6;Q\to3,4,5;R\to3,4,5;S\to1,2,6$}
\opt{B}{$\mathrm P\to1,2,6;Q\to1,2,3;R\to4,5,6;S\to4,5,6$}
\opt{C}{$\mathrm P\to1,2,3;Q\to1,2,6;R\to1,5,6;S\to3,4,5$}
\opt{D}{$\mathrm P\to1,2,6;Q\to1,2,6;R\to3,4,5;S\to3,4,5$}
\end{qoptions}
\qn{9} A very small current carrying square loop (current $I$) of side $L$ is placed in $y$-$z$ plane with centre at origin of the coordinate system (shown in figure). In column-I the coordinate of the points are given \& in column-II magnitude of strength of magnetic field is given. Then :-
\qfig{p2_phy_s2_q09_1}
\begin{center}
{\footnotesize\setlength{\tabcolsep}{2pt}\begin{tabular}{|c|p{0.3\linewidth}|c|c|}\hline
\multicolumn{2}{|c|}{List-I} & \multicolumn{2}{c|}{List-II}\\\hline
(P) & At point $O$ $(0,0,0)$ & (1) & $\dfrac{\mu_0IL^2}{2\pi a^3}$\\\hline
(Q) & At point $P_1(a,0,0)$ (here $a\gg L$) & (2) & $\dfrac{2\sqrt2\mu_0I}{\pi L}$\\\hline
(R) & At point $P_2$ $(0,a,0)$ (here $a\gg L$) & (3) & $\dfrac{\pi_0\sqrt5IL^2}{16\pi a^3}$\\\hline
(S) & At point $P_3$ $(a,a,0)$ (here $a\gg L$) & (4) & $\dfrac{\mu_0IL^2}{4\pi a^3}$\\\hline
\end{tabular}}
\end{center}
\begin{qoptions}
\opt{A}{$\mathrm P\to4;Q\to2;R\to3;S\to1$}
\opt{B}{$\mathrm P\to1;Q\to1;R\to2;S\to3$}
\opt{C}{$\mathrm P\to4;Q\to1;R\to2;S\to3$}
\opt{D}{$\mathrm P\to2;Q\to1;R\to4;S\to3$}
\end{qoptions}
\qn{10} Match the following
\begin{center}
{\footnotesize\setlength{\tabcolsep}{2pt}\begin{tabular}{|c|p{0.38\linewidth}|c|p{0.28\linewidth}|}\hline
\multicolumn{2}{|c|}{List-I} & \multicolumn{2}{c|}{List-II}\\\hline
(P) & Solid metal sphere in pure translational motion parallel to inclined plane. $B$ = uniform magnetic field directed perpendicular to the plane of motion\par\qfigi{p2_phy_s2_q09_2} & (1) & Maximum potential difference $=BrV$ between two points of sphere.\\\hline
(Q) & Solid metal sphere is in pure rolling. $B$ = uniform magnetic field directed perpendicular to the plane of motion\par\qfigi{p2_phy_s2_q09_3} & (2) & Potential difference across P and Q $=V_P-V_Q=2BrV$\\\hline
(R) & Solid metal sphere is in pure translation parallel to inclined plane in the plane of the paper\par\qfigi{p2_phy_s2_q09_4}\par $B$ = uniform horizontal magnetic field in the plane of motion & (3) & Maximum potential difference $=2BrV$ between two points of sphere.\\\hline
(S) & Solid metal is in pure rolling motion in the plane of the paper.\par $B$ = uniform magnetic field in the plane of motion.\par\qfigi{p2_phy_s2_q09_5} & (4) & Potential difference across P and Q is zero. $V_P-V_Q=0$\\\hline
 & & (5) & Potential difference across points R and S is zero i.e. $V_R-V_S=0$\\\hline
\end{tabular}}
\end{center}
\begin{qoptions}
\opt{A}{$\mathrm P\to2,3,5;Q\to2,3,5;R\to1,4,5;S\to1,4,5$}
\opt{B}{$\mathrm P\to1,2;Q\to3,4;R\to2;S\to5$}
\opt{C}{$\mathrm P\to2,3,4;Q\to1,2,3;R\to2,3;S\to1,2,3$}
\opt{D}{$\mathrm P\to2,5;Q\to3,4;R\to1,2;S\to1,5$}
\end{qoptions}
\sect{SECTION-II}
\qn{1} Two metal balls of radius $R$ are connected to each other with a thin metal wire. The balls are in a homogeneous electric field of strength $E$, The length of the metal wire is $\ell$ ; furthermore $\ell>R$. The system is at equilibrium. Find the mechanical tension $T$ in the wire. If $T$ is given by $\alpha\times10^{-12}$ N fill $\alpha$. (Given $E=\sqrt{18}$ N/C, $\ell=3$ m, $R=2$ cm)
\qn{2} A charged particle of mass $m$ and charge $q$ is shot along a horizontal rough plane with coefficient of friction $\mu$. A uniform magnetic field is applied at right angles to the plane. If the charged particle enters the magnetic field region perpendicular to the uniform field with velocity $v_0$, then the radius of curvature of the path after time $t=\dfrac{\mu_0}{8\mu g}$ is given by $r=\dfrac{7mv_0}{nqB}$, where $n$ is an integer. What is the value of $n$ ?
\qn{3} A cone rolls without slipping on a table. The half-angle at the vertex is $\alpha$, and the axis has length $h$ (see Fig.). Let the speed of the center of the base, point P in the figure, be $v$. What is the angular velocity (in rad/sec) of the cone with respect to the lab frame at the instant shown? (Given $V=20$ m/s, $h=2$m \& $\alpha=37^\circ$)
\qfig{p2_phy_s3_q03_1}
\qn{4} A fixed cone stands on its tip, with its axis in the vertical direction. The half-angle at the vertex is $\theta$. A small ring of radius $r$ rolls without slipping on the inside surface. Assume that conditions have been set up so that (1) the point of contact between the ring and the cone moves in a circle at height $h$ above the tip, and (2) the plane of the ring is at all times perpendicular to the line joining the point of contact and the tip of the cone (see Fig. ). What is the angular frequency, $\Omega$, of this circular motion? Work in the approximation where $r$ is much smaller than the radius of the circular motion, $h\tan\theta$. (Given $\theta=37^\circ$ \& $h=0.8$ m))
\qfig{p2_phy_s3_q04_1}
\qn{5} A conducting rod of length $\ell=2\sqrt5$ meter and mass $m=4$ kg lies on the horizontal table. Coefficient of friction between the rod and the table is $\mu=\left(\dfrac12\right)$. If the current in the conductor is 2 ampere, then find the minimum magnitude of magnetic field strength (in Tesla) such that conducting rod just starts to translate along $x$-axis. (take $g=10$ m/sec$^2$) [neglect the radius of rod]
\qfig{p2_phy_s3_q05_1}
\qn{6} A projectile projected tangentially from the surface of a planet of radius $R$ (64 Km) from a point A, follows the indicated trajectory during its motion. If it is at a height of $3R$ at the farthest point of its trajectory, then the velocity of projection at A, $v_0$ is given by :- (Acceleration due to gravity on surface $g=2$ m/s$^2$)
\qfig{p2_phy_s3_q06_1}
\qn{7} A bullet is fired vertically upwards with velocity $v$ from the surface of a spherical planet. When it reaches its maximum height, its acceleration due to the planet's gravity is $\dfrac14$th of its value of the surface of the planet. If the escape velocity form the planet is $v_{esc}=v\sqrt N$, then the value of $N$ is (ignore energy loss due to atmosphere)
\qn{8} Diagram shows a long wire carrying a current $i$ and a loop kept in the plane of paper. The loop is moved with some external force such that the loop reaches at a distance twice the initial. The resistance of wire used to form the loop is $\lambda\ \Omega$/m. Find the total charge (in nC) flown through the loop.\\
(Given $i=2$ amp, $a=2$ cm, $b=10$ cm, $c=5$ cm, $\lambda=2\ \Omega$/m, $\ell n\,2=0.7$, $\ell n\,3=1.1$, $\ell n\,7=1.9$)
\qfig{p2_phy_s3_q08_1}
\partsec{2}{CHEMISTRY}
\sect{SECTION-I(i)}
\qn{1} Among the following complexes (K--P)\\
$\mathrm{K_3[Fe(CN)_6]}$ \textbf{(K)}, $\mathrm{[Co(NH_3)_6]Cl_3}$ \textbf{(L)}, $\mathrm{Na_3[Co(oxalate)_3]}$ \textbf{(M)}, $\mathrm{[Ni(H_2O)_6]Cl_2}$ \textbf{(N)}, $\mathrm{K_2[Pt(CN)_4]}$ \textbf{(O)} and $\mathrm{[Zn(H_2O)_6]\,(NO_3)_2}$\textbf{(P)}\\
The diamagnetic complex are -
\begin{qoptions}
\opt{A}{K, L, N}
\opt{B}{M, O, P}
\opt{C}{L, M, O}
\opt{D}{L, M, N}
\end{qoptions}
\qn{2} Which of the following statement is/are \textbf{CORRECT} ?
\begin{qoptions}
\opt{A}{$\mathrm{[Pd^{+2}(NH_2{-}CH(CH_3){-}CO_2^-)_2]^0}$ and $\mathrm{[Zn(gly)_2]}$ complex does not show geometrical isomerism but show optical isomerism.}
\opt{B}{In $\mathrm{[PtCl_2(NCS)_2]^{-2}}$ geometrical isomerism exist, while optical isomerism does not exist}
\opt{C}{$\mathrm{Na_2[Pt(SCN)_2(ox)_2]}$ complex show structural and stereo isomerism}
\opt{D}{$\mathrm{[NiCl_2(PPh_3)_2]}$ ($\mu=\sqrt8$ B.M.) complex show geometrical isomerism}
\end{qoptions}
\qn{3} Choose correct statements :
\begin{qoptions}
\opt{A}{Good yields of R--I may be obtained by heating alcohols with sodium or potassium iodide in 95\% orthophosphoric acid.}
\opt{B}{Rate of solvolysis of 1-chloro-2-pentene is lesser than 1-chloropentane in water.}
\opt{C}{$\mathrm{NaI+H_2SO_4}$ is not a good reagent to produce alkyl iodide from alcohol because $\mathrm{H_2SO_4}$ oxidises $\mathrm{I^\ominus}$ into $\mathrm{I_2}$}
\opt{D}{Tetrahydrofuran (THF) can be synthesized by treating 4-chloro-1-butanol with aqueous sodium hydroxide.}
\end{qoptions}
\qn{4} Consider following reactions
\qfig{p2_chem_s1_q04_1}
\qfig{p2_chem_s1_q04_2}
Choose incorrect options :
\begin{qoptions}
\opt{A}{Total number of maximum atom in one plane is 8 in product B}
\opt{B}{D and F are identical}
\opt{C}{Q to E is formed by $\mathrm{S_N2}$ mechanism}
\opt{D}{Configuration is oppossite in D and F but same in D and B at chiral centre}
\end{qoptions}
\qn{5} A binary solid (AB) has rock salt structure. If the edge length of its unit cell is 400 pm, and the radius of cation is 80 pm. Select the correct statement(s) :
\begin{qoptions}
\opt{A}{Radius of anion ($\mathrm{B^-}$) is 120 pm.}
\opt{B}{Closest distance between anion - anion is $200\sqrt2$ pm.}
\opt{C}{Anion is closest surrounded by cations at a distance 200 pm.}
\opt{D}{Anions form FCC lattice with corner anion in contact with face centred anion.}
\end{qoptions}
\qn{6} Select the correct statement(s) about properties of solid.
\begin{qoptions}
\opt{A}{In antiferromagnetic substance magnetic moments of domains are oppositely oriented and cancel out each other's magnetic moment.}
\opt{B}{$\mathrm{H_2O}$, NaCl and $\mathrm{C_6H_6}$ are weakly repelled by external magnetic field.}
\opt{C}{$\mathrm{CrO_2}$ is strongly attracted by external magnetic field.}
\opt{D}{Metals conduct electricity in solid as well as in molten state.}
\end{qoptions}
\sect{SECTION-I(ii)}
\qn{7} Match the list-
\begin{center}
{\footnotesize\setlength{\tabcolsep}{2pt}\begin{tabular}{|c|p{0.42\linewidth}|c|p{0.26\linewidth}|}\hline
\multicolumn{2}{|c|}{List-I (Reactions)} & \multicolumn{2}{c|}{List-II (About major product)}\\\hline
(P) & \qfigi{p2_chem_s2_q07_1}$\xrightarrow{\mathrm{HBr}}$ & (1) & Atleast one of the product is alcohol\\\hline
(Q) & \qfigi{p2_chem_s2_q07_2}$\xrightarrow{\mathrm{HI}}$ & (2) & Atleast one product is alkyl halide\\\hline
(R) & \qfigi{p2_chem_s2_q07_3}$\xrightarrow[\text{(excess)}]{\mathrm{HI}}$ & (3) & No reaction\\\hline
(S) & \qfigi{p2_chem_s2_q07_4}$\xrightarrow[\mathrm{H_2SO_4}]{\mathrm{dil.}}$ & (4) & More than two organic product is obtained\\\hline
 & & (5) & Reaction follows nucleophilic substitution\\\hline
\end{tabular}}
\end{center}
Choose correct options
\begin{qoptions}
\opt{A}{$\mathrm P\to3,4;Q\to4,5;R\to1;S\to1,2,3$}
\opt{B}{$\mathrm P\to2,5;Q\to3;R\to2,5;S\to1,3,5$}
\opt{C}{$\mathrm P\to2,5;Q\to3;R\to2,4,5;S\to1,4,5$}
\opt{D}{$\mathrm P\to2,5;Q\to1,2,5;R\to2,4,5;S\to1,4,5$}
\end{qoptions}
\qn{8} Match the crystals with their structural characteristics :
\begin{center}
{\footnotesize\setlength{\tabcolsep}{2pt}\begin{tabular}{|c|l|c|p{0.4\linewidth}|}\hline
\multicolumn{2}{|c|}{List-I} & \multicolumn{2}{c|}{List-II}\\\hline
(P) & ZnS crystal & (1) & FCC\\\hline
(Q) & $\mathrm{CaF_2}$ crystal & (2) & Simple cubic\\\hline
(R) & NaCl crystal & (3) & Distance between closest particles is $\left(\dfrac{\sqrt3}{4}\right)a$\\\hline
(S) & Diamond crystal & (4) & Only one type of voids are occupied\\\hline
 & & (5) & Distance between closest particles is $\left(\sqrt{\dfrac14}\right)a$\\\hline
\end{tabular}}
\end{center}
\begin{qoptions}
\opt{A}{$\mathrm P\to2;Q\to3;R\to5;S\to1$}
\opt{B}{$\mathrm P\to1;Q\to5;R\to4;S\to2$}
\opt{C}{$\mathrm P\to3;Q\to4;R\to5;S\to1$}
\opt{D}{$\mathrm P\to4;Q\to5;R\to1;S\to3$}
\end{qoptions}
\qn{9} 
\begin{center}
{\footnotesize\setlength{\tabcolsep}{2pt}\begin{tabular}{|c|l|c|p{0.4\linewidth}|}\hline
\multicolumn{2}{|c|}{List-I} & \multicolumn{2}{c|}{List-II}\\\hline
(P) & Benzene + Toluene & (1) & Maximum boiling azeotrope\\\hline
(Q) & $\mathrm{H_2O+HNO_3}$ & (2) & Raoult's law is obeyed\\\hline
(R) & $\mathrm{CS_2+CH_3COCH_3}$ & (3) & Positive deviation from Raoult's law\\\hline
(S) & $\mathrm{H_2O+HCOOH}$ & (4) & Fractional distillation column gives azeotrope from residual liquid\\\hline
\end{tabular}}
\end{center}
Code :
\begin{qoptions}
\opt{A}{$\mathrm P\to2;Q\to1,4;R\to3;S\to1,4$}
\opt{B}{$\mathrm P\to2;Q\to3;R\to1;S\to4$}
\opt{C}{$\mathrm P\to1;Q\to2;R\to3;S\to4$}
\opt{D}{$\mathrm P\to1,4;Q\to3;R\to2;S\to1,3$}
\end{qoptions}
\qn{10} 
\begin{center}
{\footnotesize\setlength{\tabcolsep}{2pt}\begin{tabular}{|c|l|c|p{0.4\linewidth}|}\hline
\multicolumn{2}{|c|}{List-I} & \multicolumn{2}{c|}{List-II}\\\hline
(P) & $\mathrm{[CoCl(en)_2(NH_3)]^{2+}}$ & (1) & no optically active isomer\\\hline
(Q) & $\mathrm{[CoCl_2(en)(NH_3)_2]^{2+}}$ & (2) & number of geometrical isomers equals to number of optically active isomers\\\hline
(R) & $\mathrm{[PtCl_2(NH_3)_2(Py)_2]^{2+}}$ & (3) & Two optically active isomer\\\hline
(S) & $\mathrm{[Co(NH_3)_3(H_2O)Cl_2]^+}$ & (4) & Three geometrical isomers\\\hline
\end{tabular}}
\end{center}
\begin{qoptions}
\opt{A}{$\mathrm P\to4;Q\to3;R\to2;S\to1$}
\opt{B}{$\mathrm P\to2;Q\to4;R\to3;S\to1$}
\opt{C}{$\mathrm P\to2;Q\to4;R\to1;S\to3$}
\opt{D}{$\mathrm P\to4;Q\to3;R\to3;S\to2$}
\end{qoptions}
\sect{SECTION-II}
\qn{1} For the complex ion, $\mathrm{[Co(NH_3)_4Cl_2]^+}$\\
Total number of axial d-orbitals used for hybridisation of central atom $=x$\\
Total number of stereoisomers $=p$\\
Maximum number of atoms that can lie in same plane $=q$\\
The value of $\left(\dfrac{q}{p.x}\right)$ is -
\qn{2} Which of the following compounds/species are incorrectly matched against the indicated names.\\
(a) $\mathrm{[PdI_2(ONO)_2(H_2O)_2]}\to$ Diaquadiiodidodinitrito-N-palladium(IV)\\
(b) $\mathrm{Fe(\pi\text{-}C_5H_5)_2}\to$ Bis($\eta^5$--cyclopentadienyl)iron(II)\\
(c) $\mathrm{H_2[PtCl_6]}\to$ Hexachloridoplatinate(IV)hydride\\
(d) $\mathrm{[Al(OH)(H_2O)_5]SO_4}\to$ Pentaaquahydroxidoaluminium(III) sulphate\\
(e) $\mathrm{K_2[OsCl_5N]}\to$ Potassium nitridopentachloridoosmate(VI)\\
(f) $\mathrm{Cr(C_6H_6)_2}\to$ Bis($\eta^6$--benzene)chromium(0)
\qn{3} Calculate the total number of carbonate ores from the following in which d block metal present in (+2) oxidation state.\\
(i) Argentite \quad (ii) Calamine \quad (iii) Malachite \quad (iv) Siderite\\
(v) Galena \quad (vi) Magnetite \quad (vii) Dolomite \quad (viii) Cassiterite\\
(ix) Bauxite
\qn{4} \qfigi{p2_chem_s3_q04_1}\\
Find total number of alpha hydrogen in major product B?
\qn{5} Consider the following compounds :
\qfig{p2_chem_s3_q05_1}
Total number of aromatic compound is (R)\\
Serial number of most basic compound is (P)\\
Serial number of most acidic compound is (Q)\\
\textbf{Number of compound represented in bracket} then find value of $P+Q+R$?
\qn{6} Consider following reaction
\qfig{p2_chem_s3_q06_1}
\qfig{p2_chem_s3_q06_2}
If moles of HI consumed is $X$ and total number of Cl atom in compound (M) is $Y$ then find value of $Y-X$?
\qn{7} An unspecified quantity of an ideal gas has an initial pressure of 5 atm and temperature of $30^\circ$C. The gas is expanded at $30^\circ$C until the volume has increased by 60\% of the initial value. Next the quantity of the gas in the vessel is increased by 20\% of the initial value while the volume is maintained constant. Finally the temperature is adjusted at constant volume until the gas pressure is again 5 atm. What is the final temperature (in Kelvin) ?
\qn{8} The reaction $\mathrm{2A+B_2\to2\,AB}$ is first order with respect to $[B_2]$ and zero order with respect to $[A]$. The half-life for the overall reaction is 2 minutes when concentrations of $[A]$ and $[B]$ are in the 2 : 1 molar ratio. If 0.100 moles of A and 0.0300 moles of $\mathrm B_2$ are dissolved in 10 mL of solvent, what will concentration of $[A]$ be, 6.00 minutes after the reaction begins ?
\partsec{3}{MATHEMATICS}
\sect{SECTION-I(i)}
\qn{1} If $e^{\sin(x^2+y^2)}=\tan\dfrac{y^2}{4}+\sin^{-1}x$, then $y'(0)$ can be -
\begin{qoptions}
\opt{A}{$\dfrac{1}{3\sqrt\pi}$}
\opt{B}{$-\dfrac{1}{3\sqrt\pi}$}
\opt{C}{$-\dfrac{1}{5\sqrt\pi}$}
\opt{D}{$-\dfrac{1}{3\sqrt{5\pi}}$}
\end{qoptions}
\qn{2} If $\cot\theta+\tan\theta=x$ and $\sec\theta-\cos\theta=y$ wherever defined the which of following is/are true ?
\begin{qoptions}
\opt{A}{$x\sin\theta\cos\theta=1$}
\opt{B}{$\sin^2\theta=y\cos\theta$}
\opt{C}{$(x^2y)^{2/3}+(xy^2)^{2/3}=1$}
\opt{D}{$(x^2y)^{2/3}-(xy^2)^{2/3}=1$}
\end{qoptions}
\qn{3} For the curve $\sin x+\sin y=1$, $\displaystyle\lim_{x\to0}(x^\alpha)\dfrac{d^2y}{dx^2}$ exists then $\alpha$ can be
\begin{qoptions}
\opt{A}{$\dfrac32$}
\opt{B}{$2$}
\opt{C}{$\dfrac52$}
\opt{D}{$1$}
\end{qoptions}
\qn{4} Let $x^{\frac23}+y^{\frac23}=a^{\frac23}$, ($a>0$) be the equation of curve then which of the following is/are correct ?
\begin{qoptions}
\opt{A}{Equation of tangent at $(x_1,y_1)$ on the curve is $\dfrac{x}{x_1^{1/3}}+\dfrac{y}{y_1^{1/3}}=a^{2/3}$}
\opt{B}{Length of portion of tangent intercepted between the coordinate axes is constant.}
\opt{C}{If tangent at any point on the curve meet the coordinate axes in A and B, then locus of midpoint of AB is $x^2+y^2=\dfrac{a^2}{4}$}
\opt{D}{Equation of normal at $(x_1,y_1)$ on the curve is $xx_1^{1/3}-yy_1^{1/3}=x_1^{4/3}+y_1^{4/3}$}
\end{qoptions}
\qn{5} Consider $f(x)=(x-1)\tan^{-1}x-\dfrac12\ell n\left(1+x^2\right)$, $x\geqslant1$. Identify which of the following statement(s) is(are) correct ?
\begin{qoptions}
\opt{A}{$\displaystyle\lim_{x\to\infty}f'(x)=\dfrac\pi2$}
\opt{B}{$f'(x)$ is increasing in $[1,\infty)$}
\opt{C}{$f(x+4)-f(x)>2\pi\ \forall\ x\in[1,\infty)$}
\opt{D}{$f(x+4)-f(x)<2\pi\ \forall\ x\in[1,\infty)$}
\end{qoptions}
\qn{6} If the curves $y=2(x-a)^2$ and $y=e^{2x}$ touches each other, then $a$ is less than -
\begin{qoptions}
\opt{A}{$-1$}
\opt{B}{$0$}
\opt{C}{$1$}
\opt{D}{$2$}
\end{qoptions}
\sect{SECTION-I(ii)}
\qn{7} Let $\phi(x)=3f\left(\dfrac{x^2}{3}\right)+f\left(3-x^2\right)\ \forall\ x\in R$ where $f''(x)>0\ \forall\ x\in R$,\\
and $g(x)=(a^2-3a+2)\left(\cos^2\dfrac x4-\sin^2\dfrac x4\right)+(a-1)x+\sin1$\\
Match List-I with List-II and select the correct answer using the code given below the list.
\begin{center}
{\footnotesize\setlength{\tabcolsep}{2pt}\begin{tabular}{|c|p{0.38\linewidth}|c|p{0.3\linewidth}|}\hline
\multicolumn{2}{|c|}{List-I} & \multicolumn{2}{c|}{List-II}\\\hline
(P) & $\phi(x)$ is increasing in & (1) & $(0,1)\cup(1,4)$\\\hline
(Q) & $\phi(x)$ is decreasing in & (2) & $\left[-\dfrac32,0\right]\cup\left[\dfrac32,\infty\right)$\\\hline
(R) & The set of values of $a$ for which $g(x)$ does not possess critical point & (3) & $\left(-\infty,-\dfrac32\right]\cup\left[0,\dfrac32\right]$\\\hline
(S) & The set of values of $a$ for which $g(x)$ possess critical point is & (4) & $(-\infty,0]\cup\{1\}\cup[4,\infty)$\\\hline
\end{tabular}}
\end{center}
\begin{qoptions}
\opt{A}{$\mathrm P\to2;Q\to3;R\to4;S\to4$}
\opt{B}{$\mathrm P\to2;Q\to3;R\to4;S\to1$}
\opt{C}{$\mathrm P\to3;Q\to2;R\to1;S\to4$}
\opt{D}{$\mathrm P\to3;Q\to2;R\to4;S\to1$}
\end{qoptions}
\qn{8} Match List-I with List-II and select the correct answer using the code given below the list.
\begin{center}
{\footnotesize\setlength{\tabcolsep}{2pt}\begin{tabular}{|c|p{0.5\linewidth}|c|c|}\hline
\multicolumn{2}{|c|}{List-I} & \multicolumn{2}{c|}{List-II}\\\hline
(P) & Number of solutions of $4(\sin^6\theta+\cos^6\theta)=10-18\sin\theta\cos\theta$ in $[0,3\pi]$, is & (1) & $2$\\\hline
(Q) & Number of possible pairs $(x,y)$ satisfying $x^2-6x+10-\sin^2y=0$, where $y\in[0,2\pi]$ and $x\in R$, is & (2) & $3$\\\hline
(R) & Number of integral values of $x$ satisfying $\sin x<\dfrac12$ in $[0,2\pi]$, is & (3) & $4$\\\hline
(S) & Number of solutions satisfying $\tan x+\sec x=2\cos x$ in $[0,4\pi]$, is & (4) & $5$\\\hline
\end{tabular}}
\end{center}
\begin{qoptions}
\opt{A}{$\mathrm P\to2;Q\to1;R\to4;S\to3$}
\opt{B}{$\mathrm P\to1;Q\to2;R\to4;S\to3$}
\opt{C}{$\mathrm P\to1;Q\to2;R\to3;S\to4$}
\opt{D}{$\mathrm P\to2;Q\to1;R\to3;S\to4$}
\end{qoptions}
\qn{9} Match List-I with List-II and select the correct answer using the code given below the list.
\begin{center}
{\footnotesize\setlength{\tabcolsep}{2pt}\begin{tabular}{|c|p{0.5\linewidth}|c|c|}\hline
\multicolumn{2}{|c|}{List-I} & \multicolumn{2}{c|}{List-II}\\\hline
(P) & If $x=\alpha$ satisfies $\sin^{-1}x+\cos^{-1}x^2+\dfrac\pi2=0$, then the value of $\dfrac{\sec^{-1}\alpha-\tan^{-1}\alpha}{\cot^{-1}\alpha-\mathrm{cosec}^{-1}\alpha}$ is & (1) & $0$\\\hline
(Q) & If $a=\displaystyle\lim_{x\to\infty}\left(2.3^x+3.4^x+4.5^x\right)^{\frac1x}$ and $b=\displaystyle\lim_{x\to(-\infty)}\left(2.3^x+3.4^x+4.5^x\right)^{\frac1x}$ then $\left(\dfrac{a+b}{a-b}\right)$ is equal to & (2) & $4$\\\hline
(R) & Number of points of discontinuity of $f(x)=x\sin(\pi(2x-[x]))$ in $[0,2]$ will be\par (Note $[k]$ denotes greatest integer function less than or equal to $k$) & (3) & $1$\\\hline
(S) & $\displaystyle\lim_{x\to0}\left(\dfrac{3e^{x/3}-x-3}{x^2}\right)$ is equal to & (4) & $\dfrac13$\\\hline
 & & (5) & $\dfrac12$\\\hline
 & & (6) & $\dfrac16$\\\hline
\end{tabular}}
\end{center}
\begin{qoptions}
\opt{A}{$\mathrm P\to1;Q\to2;R\to2;S\to4$}
\opt{B}{$\mathrm P\to3;Q\to1;R\to2;S\to6$}
\opt{C}{$\mathrm P\to3;Q\to2;R\to1;S\to6$}
\opt{D}{$\mathrm P\to1;Q\to3;R\to2;S\to4$}
\end{qoptions}
\qn{10} Match List-I with List-II and select the correct answer using the code given below the list.
\begin{center}
{\footnotesize\setlength{\tabcolsep}{2pt}\begin{tabular}{|c|p{0.5\linewidth}|c|c|}\hline
\multicolumn{2}{|c|}{List-I} & \multicolumn{2}{c|}{List-II}\\\hline
(P) & If $g(x)=f(-x+xf(x))$ and $f(0)=1$, $f'(0)=2$, then value of $g'(0)$ is equal to & (1) & $\dfrac14$\\\hline
(Q) & If $e^y=e^xx^{-y}$, then $\left(\dfrac{dy}{dx}\right)_{x=e}$ is equal to & (2) & $6$\\\hline
(R) & Number of points of discontinuity of $y=[2x]$ in $[-1,2]$ is equal to (here $[.]$ denote greatest integer function) & (3) & $0$\\\hline
(S) & If $I=\displaystyle\lim_{x\to0}(1+\tan x)^{\frac1{\sin x}}$ then $[I]$ is equal to (here $[.]$ denote greatest integer function) & (4) & $2$\\\hline
 & & (5) & $3$\\\hline
\end{tabular}}
\end{center}
\begin{qoptions}
\opt{A}{$\mathrm P\to3;Q\to1;R\to2;S\to4$}
\opt{B}{$\mathrm P\to1;Q\to2;R\to3;S\to5$}
\opt{C}{$\mathrm P\to2;Q\to1;R\to4;S\to5$}
\opt{D}{$\mathrm P\to2;Q\to1;R\to3;S\to5$}
\end{qoptions}
\sect{SECTION-II}
\qn{1} If a particle starts moving from $(0,0)$ and rate of change of its abscissa is 8 cm/s on the curve $y=\sqrt x$, then rate of change of area of triangle formed by co-ordinate axes \& tangent at a point reached after 2 sec of start is
\qn{2} If $y=\left(\ell n\left(x+\sqrt{x^2+1}\right)\right)^2$ and $(x^2+1)\dfrac{d^3y}{dx^3}+(px+q)\dfrac{d^2y}{dx^2}+r\dfrac{dy}{dx}=sy$, then value of $(p+q+r+s)$ is
\qn{3} Number of points where $f(x)=\left|(2^x-1).(3^x-3).x.(x-1).(x-2)\right|$ is not differentiable
\qn{4} Let $S_n=\displaystyle\sum_{k=0}^{n}\tan^{-1}\dfrac{49}{49k^2+63k+57}$ then $\cot\left(\displaystyle\lim_{n\to\infty}S_n\right)$ is equal to
\qn{5} If minimum value of $y=\dfrac{\left(\sin^{-1}x\right)^3}{3}+\dfrac\pi2\left(\cos^{-1}x\right)^2$ exist at $x=\alpha$, then $\cos^{-1}\left(\dfrac{\cos^{-1}(-\alpha)}{\pi}\right)$ is equal to - (Use $\pi=\dfrac{22}{7}$)
\qn{6} Let $f(x)=x+\dfrac{x^2}{2}+\dfrac{x^3}{3}+\dfrac{x^4}{4}+\dfrac{x^5}{5}$ and $g(x)=f^{-1}(x)$, then $|g''(0)|$ is
\qn{7} The number of solution of equation $2\sin2x+\sin x\,(1-2\cos2x)-4=0$ in the interval $(0,\pi)$ :
\qn{8} Let $f(x)=g(x).\mathrm{sgn}(x^3+6x^2+11x+6)$, where sgn denotes signum function \& $g(x)=ax^3+x^2+x+1$, then minimum number of points of discontinuity of the function $f(x)$ is
\end{multicols}
\end{document}
